\section{相关研究工作}
\label{sec:related}

\subsection{生成式三维资产与部件结构}
三维生成模型逐渐能够从单幅图像恢复具有丰富外观的资产，但几何生成与部件理解仍是两个不同问题。TRELLIS.2 以 O-Voxel 联合表示几何与外观，并通过稀疏压缩表示进行高分辨率三维生成\citep{xiang2025trellis2}。这类表示为后续部件分析提供了几何基础，却不意味着导出的每个区域天然带有稳定的语义身份。一个外观完整的车轮仍可能与车身共用网格、材质或纹理，难以直接作为独立语义对象操作。

另一条路线将部件结构纳入生成过程。PartCrafter 采用组合式潜空间和部件内、部件间注意力，从图像联合生成可分解的三维部件\citep{lin2025partcrafter}；PartGen 则结合多视图部件分割、遮挡补全与三维重建，恢复独立部件\citep{chen2024partgen}。这些工作关注部件的生成与完整几何恢复。本文关注已有生成资产的部件划分条件：保持三维推理模型冻结，研究如何将视觉语言模型给出的部件意图转化为二维空间引导。因此，本文不以生成互不相连的完整部件、补全隐藏结构或实现可动关节为目标。

\subsection{类别无关与可提示三维部件分割}
类别无关方法试图在缺少固定类别词表时发现对象内部的结构。PartField 将二维和三维部件提议蒸馏为连续三维特征场，并通过聚类形成层次分解\citep{liu2025partfield}。P3-SAM 使用原生三维特征、多个分割头和质量预测支持点提示分割，并提供自动部件分解方式\citep{ma2025p3sam}。PartSAM 则通过三平面双分支编码器学习空间结构化特征，支持提示分割与完整对象分解\citep{zhu2025partsam}。这些方法说明，部件粒度、几何特征和提示形式共同影响输出分区；区域数量较多本身不能作为分割更准确的证据。

二维基础模型也被用于三维分割。SAM 2 提供图像与视频上的可提示掩码预测能力\citep{ravi2024sam2}，但二维掩码并不自动解决跨视图一致性和三维区域身份问题。GeoSAM2 将法线图与点图引入经几何适配的 SAM2，并通过多视图反投影与聚合获得三维部件\citep{deng2025geosam2}。SegviGen 将原生三维生成先验用于部件分割，在几何对齐的表示上以颜色表征部件身份\citep{li2026segvigen}。本文使用 SegViGen 作为冻结的二维到三维处理环节，将研究重点放在其输入引导的构建、反馈和融合，而非重新训练三维编码器或提出新的三维生成主干。

\subsection{视觉语言引导与部件语义命名}
语言可以指定所需的部件概念，但概念名称与实际像素、表面区域之间还需要空间关联。Find3D 利用二维基础模型构建部件标注数据，并通过对比学习获得支持自由文本查询的三维分割模型\citep{ma2024find3d}。SAMPart3D 学习多粒度三维部件特征，再利用多视图视觉语言模型赋予语义标签\citep{yang2024sampart3d}。Name That Part 将几何、外观与语言描述关联，通过匹配实现部件分割和命名\citep{paul2025namethatpart}。因此，本文的出发点不是宣称首次为三维部件命名，而是明确语言在既有推理流水线中的作用位置。

本文采用 Qwen3-VL\citep{bai2025qwen3vl}规划可见部件，为每个部件同时输出英文名称、边界框和正负点。名称保存在结构化记录中，框与点进入 SAM2；生成的二维区域图再作为三维分割条件。这一设计将“希望分出什么”与“提示分割模型在哪里分”结合，同时保留可追溯的部件记录。与只在分割完成后追加名称不同，语言规划会通过空间提示间接改变掩码；但名称字符串本身不作为本文 SAM2 或 SegViGen 的直接文本条件。

\subsection{部件数据与评价协议}
Objaverse 提供覆盖多类对象的三维资产资源\citep{deitke2023objaverse}，适合构造具有不同几何复杂度的研究样本。PartNeXt 进一步强调细粒度、层次化部件标注及部件中心理解任务\citep{wang2025partnext}。这些资源提示了两个评价难点：同一对象可存在多个合理粒度，且“有三维模型”不等于“有完整三维部件真值”。本文使用人工二维参考衡量可见区域的一致性，并通过三维输出的固定视角投影及多视图人工评审补充评价。我们将二维引导质量、最终三维输出的可见表面质量和名称正确性分别讨论，不将它们合并成一个未经定义的“语义准确率”。
